\documentclass[11pt,a4paper]{article}

\usepackage[utf8]{inputenc}
\usepackage[T1]{fontenc}
\usepackage[margin=1in]{geometry}
\usepackage{microtype}
\usepackage{graphicx}
\usepackage{booktabs}
\usepackage{enumitem}
\usepackage{amsmath, amssymb}
\usepackage[colorlinks=true, linkcolor=blue, urlcolor=blue]{hyperref}

\setlist[itemize]{nosep, leftmargin=*, topsep=2pt, partopsep=0pt}
\setlist[enumerate]{nosep, leftmargin=*, topsep=2pt, partopsep=0pt}

\widowpenalty=10000
\clubpenalty=10000

\title{Article Title}
\author{Author Name \\ Affiliation}
\date{\today}

\begin{document}

\maketitle

\begin{abstract}
A short summary of the article's purpose, approach, and findings.
\end{abstract}

\section{Introduction}

Open with the problem and why it matters. Establish what the reader will get from this article and how it relates to prior work.

\section{Background}

Give the minimum context needed to follow the rest. Cite related work; do not summarize an entire field.

\section{Method}

Describe the approach. Equations:

\begin{equation}
  f(x) = \int_{0}^{x} g(t) \, dt
  \label{eq:integral}
\end{equation}

Equation~\ref{eq:integral} defines the integral form used below.

\section{Results}

Present findings concretely. Where appropriate, use tables:

\begin{table}[htbp]
  \centering
  \caption{Headline result.}
  \label{tab:result}
  \begin{tabular}{lrr}
    \toprule
    Method & Score & Time (s) \\
    \midrule
    Baseline & 72.1 & 1.4 \\
    Ours     & 84.6 & 1.6 \\
    \bottomrule
  \end{tabular}
\end{table}

\section{Discussion}

Interpret. State limitations explicitly.

\section{Conclusion}

One paragraph. Tie back to the problem opened in \S1.

\end{document}
